\section{Conclusion}
This study conducted a comprehensive game-theoretic analysis to evaluate the rationality and effectiveness of adopting a negotiation workflow within Large Language Models (LLMs) across a spectrum of classic strategic scenarios. By modeling interactions through well-established complete-information games, including the Prisoner's Dilemma, Stag Hunt, Battle of the Sexes, Wait-Go Game, Duopolistic Competition, Escalation Game, Monopoly Game, Draco and Harry Game, we assessed how different LLMs, specifically Claude-3.5 Sonnet, GPT-4o, and Claude-3 Opus, navigate the balance between cooperation and competition.

Expanding our investigation to more realistic settings, we explored the Deal-No-Deal Game, which incorporates incomplete-information, to assess whether LLMs can efficiently allocate resources and negotiate agreements under conditions of uncertainty. In this context, we designed a workflow based on Bayesian updates to achieve pareto optimal and envy free allocations, thereby enhancing the negotiation process.


Our findings indicate that the adoption of the workflow generally promotes pareto optimal outcomes, wherein both agents achieve higher collective payoffs compared to non-cooperative strategies. For instance, in scenarios like the Stag Hunt and Battle of the Sexes, the workflow facilitated effective negotiation and coordination, enabling LLMs to reach mutually beneficial equilibria. Similarly, in the Deal-No-Deal Game, the structured workflow enhanced decision-making processes, leading to more efficient resource allocations.
